% Provide an overview of your research area. Ensure you have identified the
% challenges in the existing solutions and the motivation for your work.
% Motivation: Clearly state the motivation behind the work. Challenges: Outline
% the challenges in the existing solution.

% =========================================================================== %
% --- Theory ---

Recent developments in the world of Artificial Intelligence and its
ability to generate almost undetectable images have taken over the
world by storm, with every field relying on technology to make
content easily and within minutes. When it comes to image generation,
however, AI tools today can fabricate convincing forgeries for
disinformation, ``deepfakes'', and fraud just as easily. Hence, the
question, where do we draw the line on synthetic content?
Traditionally, digital forensics has relied on detecting copy-move
duplication, splicing composites and object removal, left spatial
inconsistencies, using CNN architectures, but now diffusion models
like Stable Diffusion and DALL·E produce globally synthetic images
indistinguishable to the human eye which poses a problem.

On March 23, 2023, certain images were leaked of American President
Donald Trump and Russian President Vladimir Putin getting tackled by
police officers \cite{klepper2023deepfake}. While the images become
widely sensational, they were AI generated and such situations invoke
mistrust in people unable to identify what is real and what is not.
Along with image generation, AI-powered inpainting, which is the
selective modification of authentic photographs, has become
accessible to all and hugely popular, with Adobe Photoshop's
Generative Fill and Google's Gemini and many other such tools.
The consequence of undetected image forgeries is a matter of serious
concern, where political misinformation, manipulated imagery in civil
disputes and even faked identity documents can cause heavy harm and
damage to reputation.

Existing detection models, even state-of-the-art systems like
FatFormer, GRAMNet, and DiffCoR, face two critical limitations: they
generalize poorly to generator models not seen during training and
several require significant computational resources unavailable in
resource-constrained environments. A computationally efficient,
generalizable, and inpainting-aware detector is therefore a pressing
requirement.

% =========================================================================== %
% --- Problem Statement ---

AI-generated and AI-manipulated (inpainted) images that are produced
by diffusion-based generative models must be reliably detected by a
system with good accuracy and domain generalization capability. The
system must be able to classify images into three broad categories:
not AI-generated, AI-generated, or AI-manipulated. The detection must
be robust to image post-processing operations like compression or
resizing, and able to identify AI-generated images across different
generation architectures.

Current AI image generators are too powerful, and people cannot
distinguish between what is real and what has been generated. There
is a need to make detection automatic and algorithm-driven in this
scenario. Failure to detect such forgeries will have bearing
consequences in the fields of news and journalism, legal junctures,
fraud and crime prevention, social media integrity and many other
relevant spheres of life. The gap in the detection of inpainted
images is especially worrisome as bad actors can manipulate select
regions of an image to their liking while also being able to pass
them as credible against traditional verification algorithms.
